\section{结论}
本文导出了一个公式，它把杨--米尔斯梯度流产生的某些规范不变局域乘积的短流
时间行为，与杨--米尔斯理论中正确归一化的守恒能量--动量张量联系起来。我们
的主要结果是式~\eqref{eq:(4.29)}。式~\eqref{eq:(4.29)} 的右端可以在格点规
范理论中用 Wilson 流、配合算符~\eqref{eq:(3.5)} 与~\eqref{eq:(3.6)} 的适当
离散化来计算（例如见文献~\cite{Luscher:2010iy,Borsanyi:2012zs}）。这里必须
首先取连续极限~$a\to0$，然后才取 $t\to0$ 极限；否则我们的基本推理不成立。

虽然公式~\eqref{eq:(4.29)} 在数学上应当是正确的，但它的实际有用性则是另一
个问题，必须通过数值仔细检验。\footnote{我们希望在不远的将来回到这一问
题。}由于要使我们的推理有效，格距~$a$ 必须远小于流时间的平方根%
~$\sqrt{8t}$，式~\eqref{eq:(4.29)} 的可靠应用将需要相当小的格距。人们还会担
心高维算符的污染（即式~\eqref{eq:(3.7)} 与~\eqref{eq:(3.8)} 中的 $O(t)$
项）以及本文未予考虑的有限体积效应。如果我们的策略在实践中可行，它将提供
一种全新的方法来计算包含良定义能量--动量张量的关联函数。显然，本文处理格
点上能量--动量张量的思路并不限于纯杨--米尔斯理论，尽管存在其他场时处理可
能稍微复杂一些。应用将包括：切变与体黏滞系数的确定（例如见文
献~\cite{Meyer:2007ic,Meyer:2007dy}）、热力学量的测量（见文
献~\cite{Giusti:2012yj} 及其所引文献）、与（近似）伸缩不变性相联系的赝
Nambu--Goldstone 玻色子的质量与衰变常数（见文献~\cite{Appelquist:2010gy}
及其所引文献），等等。

同样显然的是，我们的基本想法——用格点正规化定义的算符可以和连续理论中的
算符通过梯度流联系起来——并不限于能量--动量张量。例如，或许可以从局域乘
积的小流时间极限出发，在格点上构造理想的手征流或理想超流。追求这一想法将
是有意义的。
